\documentclass[12pt]{article}

\input{../../_assets/preambulo.tex}

\usepackage{algorithm}
\floatname{algorithm}{Algoritmo}
\usepackage{algpseudocode}

\algrenewcommand\algorithmicif{\textbf{Si}}
\algrenewcommand\algorithmicthen{\textbf{entonces}}
\algrenewcommand\algorithmicelse{\textbf{Si no}}
\algrenewcommand\algorithmicend{\textbf{Fin}}
\algrenewtext{For}[1]{\textbf{Para} #1 \textbf{hacer}}
\algrenewtext{EndFor}{\textbf{Fin Para}}
\algrenewtext{While}[1]{\textbf{Mientras} #1 \textbf{hacer}}
\algrenewtext{EndWhile}{\textbf{Fin Mientras}}
\algrenewtext{ElsIf}[1]{\textbf{Si no, si} #1 \textbf{entonces}}

\begin{document}

    % 1. Foto de fondo
    % 2. Título
    % 3. Encabezado Izquierdo
    % 4. Color de fondo
    % 5. Coord x del titulo
    % 6. Coord y del titulo
    % 7. Fecha

    
    \input{../../_assets/portada}
    \portadaExamen{ffccA4.jpg}{Modelos \\ Avanzados \\ de Computación \\Examen III}{Modelos Avanzados de Computación. Examen III}{MidnightBlue}{-8}{28}{2026}{José Manuel Sánchez Varbas}

    \begin{description}
        \item[Asignatura] Modelos Avanzados de Computación.
        \item[Curso Académico] 2022-23.
        \item[Grado] Grado en Ingeniería Informática.
        \item[Descripción] Ordinaria.
        \item[Fecha] 12 de Junio de 2023.  
    \end{description}
    \newpage


    % ------------------------------------

    \begin{ejercicio}[2 puntos]
        Estudiar si son decidibles y/o semidecidibles los siguientes problemas:
        \begin{enumerate}
            \item Problema de las correspondencias de Post cuando el alfabeto es $A = \{1\}$.
            \item Dada una MT $M$, ¿es cierto que la complejidad del lenguaje aceptado es polinómica?
            \item Dado un programa Post--Turing $P$, ¿es cierto que $P$ no acepta la palabra vacía?
            \item Dadas dos MT $M_1$ y $M_2$, ¿es cierto que $<M_1>$ es una subcadena de $<M_2>$?
        \end{enumerate}
    \end{ejercicio}

    \begin{ejercicio}[2 puntos]
        Razona si son verdaderas o falsas las siguientes afirmaciones:
        \begin{enumerate}
            \item ¿Es cierto que si un problema tiene un algoritmo $\delta$-aproximado polinómico con $\delta = 2$, entonces también tendrá un algoritmo $\delta$-aproximado polinómico con $\delta = 3$?
            
            \item Si un problema está en \textbf{CoNP} si y solo si existe una MT no determinista tal que si la respuesta del problema es afirmativa, la MT siempre acepta y si es negativa, entonces puede aceptar o rechazar, pero al menos uno de los cálculos posibles acaba en rechazo.
            
            \item Si $P_1 \propto P_2 \propto P_3$ y $P_3$ no es semidecidible, entonces $P_1$ tampoco lo es.
            
            \item Un problema \textbf{NP}-difícil es siempre \textbf{NP}-completo.
        \end{enumerate}
    \end{ejercicio}

    \begin{ejercicio}[1.5 puntos]
        Describe la relación existente entre el problema 2-SAT y la búsqueda de caminos en grafos dirigidos.
    \end{ejercicio}

    \begin{ejercicio}[1.5 puntos]
        Define las clases de complejidad \textbf{P} y \textbf{EXP}, ¿qué relación existe entre \textbf{P} y \textbf{EXP}?
    \end{ejercicio}

    \begin{ejercicio}[1.5 puntos]
        ¿Cómo se reduce el problema del conjunto independiente al problema del cubrimiento por vértices en grafos? (versiones de decisión de ambos problemas).
    \end{ejercicio}

    \begin{ejercicio}[1.5 puntos]
        Explica cómo se puede simular una MT multicinta mediante una MT con una sola cinta. Si para una entrada de longitud $n$, la MT multicinta da como mucho $t(n)$ pasos, ¿qué podemos asegurar sobre el número de pasos de la MT con una solo cinta?
    \end{ejercicio}

    \newpage
    \setcounter{ejercicio}{0} % Reiniciar contador de ejercicios
    \noindent
    \textbf{Solución.}

    \begin{ejercicio}[2 puntos]
        Estudiar si son decidibles y/o semidecidibles los siguientes problemas:
        \begin{enumerate}
            \item Problema de las correspondencias de Post cuando el alfabeto es $A = \{1\}$. \\
            
            Tenemos que $A^{*} = \{1^n : n \geq 0\}$. Cada bloque del PCP será de la forma 
            $$\dfrac{1^{p_i}}{1^{q_i}}$$ 
            con $p_i$ la longitud de la palabra superior, y $q_i$ la longitud de la palabra inferior, $i = 1, \ldots, k$, y 
            $k$ el número de bloques. Entonces, para una sucesión no vacía de enteros $i_1, \ldots, i_m$, la palabra de arriba 
            será $1^{p_{i_1}} \cdots 1^{p_{i_m}} = 1^{p_{i_1} + \cdots + p_{i_m}}$. Análogamente, la palabra de abajo será 
            $1^{q_{i_1} + \cdots + q_{i_m}}$. Entonces, como solo hay un símbolo, dos palabras son iguales si y solo si tienen
            la misma longitud, es decir, la sucesión es válida (como solución al PCP) si y solo si 
            $p_{i_1} + \cdots + p_{i_m} = q_{i_1} + \cdots + q_{i_m}$. Definimos $d_i = p_i - q_i$. La condición de solución
            se reescribe como
            $$p_{i_1} + \cdots + p_{i_m} = q_{i_1} + \cdots + q_{i_m} \iff 
            (p_{i_1} - q_{i_1}) + \cdots + (p_{i_m} - q_{i_m}) = 0 \iff$$$$ d_{i_1} + \cdots d_{i_m} = 0$$
            Para ver que es decidible usamos el siguiente algoritmo:

            \begin{algorithm}
                \caption{MT $M_{P_1}$}
                \begin{algorithmic}
                    \State \textbf{Entrada:} Una instancia de PCP sobre $A=\{1\}$ con bloques
                    $\frac{1^{p_1}}{1^{q_1}},\dots,\frac{1^{p_k}}{1^{q_k}}$.

                    \For{$i=1,\dots,k$}
                        \State Calculamos $d_i = p_i-q_i$.
                    \EndFor

                    \If{existe $i$ tal que $d_i=0$}
                        \State Aceptamos.
                    \EndIf

                    \If{existen $i,j$ tales que $d_i>0$ y $d_j<0$}
                        \State Aceptamos.
                    \EndIf

                    \State Rechazamos.
                \end{algorithmic}
            \end{algorithm}

            Si algún $d_i = 0$, entonces el bloque $i$ es solución. Si hay $d_i > 0$ y $d_j < 0$ para $i,j \in \{1, \ldots, k\}$,
            escribimos $d_i = r > 0$, y $d_j = -s < 0$. Tomando $s$ copias del bloque $i$ y $r$ copias del bloque $j$,
            la diferencia total será $s \cdot r + r \cdot (-s) = 0$, luego hay solución. Si todos los $d_i$ son positivos,
            cualquier suma no vacía será positiva. Si, por otro lado, todos los $d_i$ son negativos, cualquier suma no vacía 
            será negativa. En ninguno de estos dos casos puede haber solución. Como el algoritmo solo recorre un número
            finito de bloques y siempre termina, el PCP sobre un alfabeto con un único símbolo es decidible.

            \item Dada una MT $M$, ¿es cierto que la complejidad del lenguaje aceptado es polinómica? \\
            
            Interpretamos el problema como
            $$
            \text{POL}(M)=\{\langle M\rangle : L(M)\in \mathbf{P}\}.
            $$

            No es decidible por el Teorema de Rice, ya que ``$L(M)\in \mathbf{P}$'' depende solo del lenguaje aceptado por $M$, no de la máquina concreta, y 
            es una propiedad no trivial de los lenguajes r.e.: por ejemplo, $\emptyset \in \mathbf{P}$ (basta considerar una MT que rechace cualquier entrada), pero el 
            lenguaje universal $L_u$ es r.e. y no decidible, luego no está en $\mathbf{P}$ (esto último se debe a que si $L_u \in \textbf{P}$, entonces existiría una MT 
            determinista que decide $L_u$ en tiempo polinómico. En particular, nunca ciclaría, luego $L_u$ sería un lenguaje recursivo, es decir, el problema UNIVERSAL 
            sería decidible, lo cual sabemos que no es cierto, por lo que $L_u \notin \textbf{P}$). \\

            Veamos además que no es semidecidible. Reducimos desde
            $$
            \text{DIAGONAL}(M)=\{\langle M\rangle : M \text{ no acepta } \langle M\rangle\},
            $$
            que no es semidecidible. Dada una MT $M$, construimos una MT $M_{POL}$:

            \begin{algorithm}
            \caption{MT $M_{POL}$}
            \begin{algorithmic}
            \State \textbf{Entrada:} una palabra $x \in \{0,1\}^{*}$
            \State Simular $M$ con entrada $\langle M\rangle$.
            \If{$M$ acepta $\langle M\rangle$}
                \State Simular una MT fija $M'$ tal que $L(M')=L_u$ con entrada $x$.
                \If{$M'$ acepta $x$}
                    \State Aceptar.
                \EndIf
            \EndIf
            \end{algorithmic}
            \end{algorithm}

            Si $\langle M\rangle \in \text{DIAGONAL}$, entonces $M$ no acepta $\langle M\rangle$, luego $L(M_{POL})=\emptyset \in \mathbf{P}$.

            Si $\langle M\rangle \notin \text{DIAGONAL}$, entonces $M$ acepta $\langle M\rangle$, luego 
            $M_{POL}$ se comporta como $M'$ sobre toda entrada $x$, y por tanto $L(M_{POL})=L(M')=L_u$, que no es decidible y por tanto no pertenece a $\mathbf{P}$.

            Así,
            $$
            \langle M\rangle \in \text{DIAGONAL}
            \Longleftrightarrow
            \langle M_{POL}\rangle \in \text{POL}.
            $$

            Como $\text{DIAGONAL} \propto \text{POL}$ y DIAGONAL no es semidecidible, $\text{POL}$ tampoco es semidecidible. Por tanto, el problema es no semidecidible y, en particular, no decidible.

            \newpage

            \item Dado un programa Post--Turing $P$, ¿es cierto que $P$ no acepta la palabra vacía? \\
            
            Definimos
            $$
            \text{NOACEPTA}_{\varepsilon}(P)=\{\langle P\rangle : P \text{ no acepta } \varepsilon\}.
            $$

            Este problema es no semidecidible. Reducimos desde
            $$
            \text{DIAGONAL}(M)=\{\langle M\rangle : M \text{ no acepta } \langle M\rangle\}.
            $$

            Dada una MT $M$, construimos un programa Post--Turing $P_M$ equivalente a la siguiente máquina:

            \begin{algorithm}
            \caption{Programa $P_M$}
            \begin{algorithmic}
            \State \textbf{Entrada:} una palabra $x$
            \State Ignorar $x$.
            \State Simular $M$ con entrada $\langle M\rangle$.
            \If{$M$ acepta $\langle M\rangle$}
                \State Aceptar.
            \EndIf
            \end{algorithmic}
            \end{algorithm}

            Esta construcción es válida porque los programas Post--Turing son equivalentes a las MT.

            Entonces:
            \begin{itemize}
                \item si $\langle M\rangle \in \text{DIAGONAL}$, entonces $M$ no acepta $\langle M\rangle$, y por construcción $P_M$ no acepta ninguna palabra; en particular, no acepta $\varepsilon$.
                \item Si $\langle M\rangle \notin \text{DIAGONAL}$, entonces $M$ acepta $\langle M\rangle$, y por construcción $P_M$ acepta todas las palabras; en particular, acepta $\varepsilon$.
            \end{itemize}

            Por tanto,
            $$
            \langle M\rangle \in \text{DIAGONAL}(M)
            \Longleftrightarrow
            \langle P_M\rangle \in \text{NOACEPTA}_{\varepsilon}(P_M).
            $$

            Como $\text{DIAGONAL} \propto \text{NOACEPTA}_\varepsilon$ y DIAGONAL no es semidecidible, $\text{NOACEPTA}_{\varepsilon}$ tampoco es semidecidible. En particular, no es decidible.

            \newpage

            \item Dadas dos MT $M_1$ y $M_2$, ¿es cierto que $<M_1>$ es una subcadena de $<M_2>$? \\
            
            Es decidible, luego también semidecidible. El problema solo depende de las codificaciones $\langle M_1\rangle$ y
            $\langle M_2\rangle$, que son dos palabras finitas, no del comportamiento de las máquinas. \\

            Basta aplicar el siguiente algoritmo:

            \begin{algorithm}
            \caption{$M_{P_4}$}
            \begin{algorithmic}
            \State \textbf{Entrada:} $\langle M_1,M_2\rangle$
            \State Comprobar si $\langle M_1\rangle$ aparece como subcadena de $\langle M_2\rangle$.
            \If{aparece}
                \State Aceptar.
            \Else
                \State Rechazar.
            \EndIf
            \end{algorithmic}
            \end{algorithm}

            Como la comprobación de subcadena entre dos palabras finitas termina siempre (como mucho se recorre toda
            la cadena de $\langle M_2 \rangle$, que es la más larga a priori, si como subcadena $\langle M_1 \rangle$ 
            se propone la propia cadena completa $\langle M_2 \rangle$), el problema es decidible.
        \end{enumerate}
    \end{ejercicio}

    \begin{ejercicio}[2 puntos]
        Razona si son verdaderas o falsas las siguientes afirmaciones:
        \begin{enumerate}
            \item ¿Es cierto que si un problema tiene un algoritmo $\delta$-aproximado polinómico con $\delta = 2$, entonces también 
            tendrá un algoritmo $\delta$-aproximado polinómico con $\delta = 3$? \\

            $\boxed{\text{Verdadera}}$. En ambos casos, minimización y maximización, la razón de
            eficacia se define de forma que valores más cercanos a $1$ son mejores. Si el
            problema es de maximización, ser $2$-aproximado significa que, para toda entrada
            $x$,
            $$
            \frac{OPT(x)}{B(ALG(x))}\leq 2.
            $$
            En particular,
            $$
            \frac{OPT(x)}{B(ALG(x))}\leq 3.
            $$
            Si el problema es de minimización, ser $2$-aproximado significa que, para toda
            entrada $x$,
            $$
            \frac{C(ALG(x))}{OPT(x)}\leq 2,
            $$
            y por tanto también
            $$
            \frac{C(ALG(x))}{OPT(x)}\leq 3.
            $$
            Así, en cualquiera de los dos casos, el mismo algoritmo polinómico es también
            $3$-aproximado.
            
            \item Si un problema está en \textbf{CoNP} si y solo si existe una MT no determinista tal que si la respuesta del problema es 
            afirmativa, la MT siempre acepta y si es negativa, entonces puede aceptar o rechazar, pero al menos uno de los cálculos 
            posibles acaba en rechazo. \\
            
            $\boxed{\text{Verdadera}}$. Un problema está en \textbf{CoNP}, por definición, si su complementario está en \textbf{NP}. Equivalentemente,
            existe una MTND polinómica tal que 
            $$x \in L \iff \text{ todas las ramas aceptan.}$$
            Por lo tanto, si $x \notin L$, no todas las ramas aceptan, es decir, al menos una rechaza.

            \item Si $P_1 \propto P_2 \propto P_3$ y $P_3$ no es semidecidible, entonces $P_1$ tampoco lo es. \\
            
            $\boxed{\text{Falsa}}$. Las cosas negativas (no decidible o no semidecidible) que podamos decir sobre los problemas 
            aplican ``de izquierda a derecha'' según el sentido de las reducciones que hemos visto en teoría. Asimismo, las cosas 
            positivas (decidible o semidecidible) que podamos decir sobre los problemas aplican ``de derecha a izquierda''
            según este mismo sentido. \\
            
            Un contrajemplo sería el que sigue. Consideramos $P_1 = \emptyset$, que es decidible, y $P_2 = P_3 = \text{DIAGONAL}$,
            que no es semidecidible. Entonces $P_1 \propto P_2 \propto P_3$, pero $P_1$ sí es semidecidible, y $P_3$ no.
            
            \item Un problema \textbf{NP}-difícil es siempre \textbf{NP}-completo. \\
            
            $\boxed{\text{Falsa}}$. Para ser \textbf{NP}-completo hacen falta dos condiciones; que el problema, pongamos $P \in \textbf{NP}$
            y que $P$ sea \textbf{NP}-difícil. Ser \textbf{NP}-difícil solo garantiza la segunda, pero no garantiza que $P \in \textbf{NP}$.
            Como contraejemplo, consideramos $\text{PARADA}$. Es NP-difícil, pues
            $SAT \propto \text{PARADA}$: dada una fórmula $\varphi$, construimos una MT
            $M_\varphi$ que prueba todas las asignaciones de verdad de $\varphi$ y para
            si encuentra una que satisface $\varphi$; si no existe ninguna, entra en un
            bucle infinito. Entonces
            $$
            \varphi \in SAT \Longleftrightarrow \langle M_\varphi,\varepsilon\rangle
            \in \text{PARADA}.
            $$
            Sin embargo, $\text{PARADA}$ no es decidible, luego no pertenece a
            $\textbf{NP}$. Por tanto, es NP-difícil pero no NP-completo.
        \end{enumerate}
    \end{ejercicio}

    \begin{ejercicio}[1.5 puntos]
        Describe la relación existente entre el problema 2-SAT y la búsqueda de caminos en grafos dirigidos. \\

        El problema $2$-SAT se reduce a estudiar caminos en un grafo dirigido, llamado grafo de implicaciones. \\

        Dada una fórmula de $2$-SAT, a cada cláusula $(a\vee b)$ le asociamos las implicaciones
        $$
        \neg a \to b,
        \qquad
        \neg b \to a.
        $$
        Construimos el grafo dirigido cuyos vértices son los literales y cuyas aristas son esas implicaciones. \\

        Entonces la fórmula es insatisfacible si y solo si existe una variable $x_i$ tal que hay camino dirigido de $x_i$ a $\neg x_i$ y también camino dirigido de $\neg x_i$ a $x_i$:
        $$
        x_i \leadsto \neg x_i
        \qquad \text{y} \qquad
        \neg x_i \leadsto x_i.
        $$

        Por tanto, $2$-SAT se decide mediante búsquedas de caminos en grafos dirigidos. Como la existencia de caminos en grafos dirigidos se resuelve en tiempo polinómico, se concluye que
        $$
        2\text{-SAT}\in \mathbf{P}.
        $$
    \end{ejercicio}

    \begin{ejercicio}[1.5 puntos]
        Define las clases de complejidad \textbf{P} y \textbf{EXP}, ¿qué relación existe entre \textbf{P} y \textbf{EXP}? \\

        La clase $\textbf{P}$ es la clase de los lenguajes decidibles en tiempo polinómico determinista:
        $$
        \textbf{P}=\bigcup_{j>0}\textbf{TIEMPO}(n^j).
        $$

        La clase $\textbf{EXP}$ es la clase de los lenguajes decidibles en tiempo exponencial determinista:
        $$
        \textbf{EXP}=\bigcup_{j>0}\textbf{TIEMPO}(2^{n^j}).
        $$

        Como todo polinomio está acotado por una exponencial, se tiene
        $$
        \textbf{P}\subseteq \textbf{EXP}.
        $$

        Además, por el Teorema de la Jerarquía en Tiempo, si $f(n)\geq n$ es una función de complejidad propia, entonces
        $$
        \textbf{TIEMPO}(f(n)) \subsetneq \textbf{TIEMPO}(f^2(n)).
        $$

        Tomando
        $$
        f(n)=2^n,
        $$
        se obtiene
        $$
        \textbf{TIEMPO}(2^n)\subsetneq \textbf{TIEMPO}((2^n)^2).
        $$

        Como
        $$
        \textbf{P}\subseteq \textbf{TIEMPO}(2^n)
        $$
        y
        $$
        \textbf{TIEMPO}((2^n)^2)\subseteq \textbf{EXP},
        $$
        existe un lenguaje en $\textbf{EXP}$ que no está en $\textbf{P}$. Por tanto,
        $$
        \textbf{P}\neq \textbf{EXP}.
        $$

        Junto con $\textbf{P}\subseteq \textbf{EXP}$, concluimos que
        $$
        \textbf{P}\subsetneq \textbf{EXP}.
        $$
    \end{ejercicio}

    \newpage

    \begin{ejercicio}[1.5 puntos]
        ¿Cómo se reduce el problema del conjunto independiente al problema del cubrimiento por vértices en grafos? 
        (versiones de decisión de ambos problemas). \\

        Consideramos las versiones de decisión:
        $$
        \text{CI}=\{(G,k): G \text{ tiene un conjunto independiente de tamaño } \geq k\},
        $$
        $$
        \text{CV}=\{(G,k): G \text{ tiene un cubrimiento por vértices de tamaño } \leq k\}.
        $$

        La reducción es:
        $$
        (G,k)\longmapsto (G, |V|-k).
        $$

        Veamos la equivalencia. Sea $I\subseteq V$. Entonces $I$ es independiente si y solo si ninguna arista tiene sus dos extremos en $I$, lo cual equivale a que toda arista tenga al menos un extremo en $V\setminus I$. Por tanto,
        $$
        I \text{ es independiente}
        \Longleftrightarrow
        V\setminus I \text{ es cubrimiento por vértices}.
        $$

        Así,
        $$
        |I|\geq k
        \Longleftrightarrow
        |V\setminus I|\leq |V|-k.
        $$

        Luego
        $$
        (G,k)\in \text{CI}
        \Longleftrightarrow
        (G,|V|-k)\in \text{CV}.
        $$

        Por tanto,
        $$
        \text{CI}\propto \text{CV}.
        $$
    \end{ejercicio}

    \begin{ejercicio}[1.5 puntos]
        Explica cómo se puede simular una MT multicinta mediante una MT con una sola cinta. Si para una entrada de longitud $n$, 
        la MT multicinta da como mucho $t(n)$ pasos, ¿qué podemos asegurar sobre el número de pasos de la MT con una solo cinta? \\

        Sea $M$ una MT con $k$ cintas. Construimos una MT $N$ con una sola cinta usando $2k$ pistas: para cada cinta de $M$, una pista 
        guarda su contenido y otra pista marca con un símbolo especial $*$ la posición del cabezal correspondiente. \\

        Para simular un paso de $M$, la máquina $N$ recorre su cinta de izquierda a derecha, localiza las $k$ marcas $*$ y guarda en 
        su control finito los $k$ símbolos leídos. Con esa información aplica la transición de $M$ y vuelve a recorrer la cinta para 
        escribir los nuevos símbolos y desplazar las marcas de los cabezales. \\

        Si $M$ da como mucho $t(n)$ pasos, entonces en cada cinta solo puede haber usado $O(t(n))$ casillas. Por tanto, cada paso de $M$ se simula recorriendo una zona de longitud $O(t(n))$. Luego el tiempo total de $N$ es
        $$
        O(t(n)\cdot t(n))=O(t^2(n)).
        $$

        Así, toda MT multicinta que trabaja en tiempo $t(n)$ puede simularse mediante una MT de una sola cinta en tiempo
        $$
        O(t^2(n)).
        $$
    \end{ejercicio}

\end{document}
